\documentclass[11pt,a4paper]{article}

\usepackage[margin=1.15in]{geometry}
\usepackage{amsmath,amssymb,amsthm,mathtools}
\usepackage{microtype}
\usepackage{hyperref}

\hypersetup{
  colorlinks=true,
  linkcolor=blue,
  citecolor=blue,
  urlcolor=blue
}

\theoremstyle{plain}
\newtheorem{theorem}{Theorem}[section]
\newtheorem{proposition}[theorem]{Proposition}
\newtheorem{lemma}[theorem]{Lemma}
\theoremstyle{definition}
\newtheorem{definition}[theorem]{Definition}
\newtheorem{postulate}{Postulate}
\theoremstyle{remark}
\newtheorem{remark}[theorem]{Remark}

\newcommand{\N}{\mathbb{N}}
\newcommand{\Z}{\mathbb{Z}}
\newcommand{\Q}{\mathbb{Q}}
\newcommand{\R}{\mathbb{R}}
\newcommand{\C}{\mathbb{C}}
\newcommand{\alef}{\aleph_{0}}
\newcommand{\cont}{\mathfrak{c}}
\newcommand{\Hil}{\mathcal{H}}
\newcommand{\Uni}{\mathcal{U}}
\newcommand{\Geo}{\mathcal{G}}
\newcommand{\Samp}{\mathsf{S}}
\newcommand{\KG}{K_{G}}
\newcommand{\UTM}{\mathsf{U}_{G}}
\newcommand{\abs}[1]{\lvert#1\rvert}
\newcommand{\norm}[1]{\lVert#1\rVert}
\newcommand{\ket}[1]{\lvert#1\rangle}
\newcommand{\bra}[1]{\langle#1\rvert}
\newcommand{\braket}[2]{\langle#1|#2\rangle}
\newcommand{\dd}{\mathrm{d}}
\newcommand{\tr}{\operatorname{tr}}
\newcommand{\Diff}{\operatorname{Diff}}

\title{\textbf{k}: A Theory of Everything\\[0.4em]
\large Quantum Field Theory as $\alef$, General Relativity as $\cont$,\\
and the Kolmogorov Merge of Sampled Universes}
\author{k}
\date{2026}

\begin{document}
\maketitle

\begin{abstract}
A theory of everything must say how the two successful infinities of twentieth-century physics communicate.
Quantum field theory is the physics of a separable Hilbert space: a countable orthonormal basis, occupation numbers, Fock modes, programs.
That is Cantor's infinity of the integers, $\alef$.
General relativity is the physics of a smooth Lorentzian $4$-manifold: uncountably many points, a metric of cardinality $\cont = 2^{\alef}$.
That is Cantor's infinity of the reals.
The continuum is the power set of the integers, so a classical universe is a completion of a countable quantum description.
We take this as ontology, not metaphor.

Universes are sampled from the quantum state of the universe.
Each sample is a finite string---Cauchy data, a code, a program---drawn from the integer skeleton of that state.
General relativity inflates the string to a continuum spacetime.
The weight of the resulting universe is the gravitational Kolmogorov complexity of the sample: Einstein's equation is the shortest program that extends $3$-data to a $4$-geometry, and $2^{-\KG}$ is the corresponding Solomonoff measure.
The physical world is the merge of these universes against those weights.
The merge is a countable sum, hence well-defined; the continuum appears only after inflation.
Classical Einstein gravity is the typical merge of low-complexity samples.
Quantum field theory is the residual integer dynamics on the merged background.
The gravitational path integral, the Born rule, black-hole entropy, the holographic bound, and the cosmological arrow of time are the same object viewed at different resolutions.
\end{abstract}

\tableofcontents
\newpage

\section{The problem of two infinities}

Physics has two working theories and two working infinities, and they are not the same infinity.

Quantum field theory (QFT) is built on a separable Hilbert space $\Hil$.
Separability means there is a countable orthonormal basis $\{\ket{n}\}_{n\in\N}$.
Fock space is assembled from occupation numbers in $\N$.
A lattice regularization, a mode expansion, a particle species, a spin component: all of these are indexed by integers.
Even $L^{2}(\R^{n})$, whose classical configuration space is a continuum, admits a countable Hilbert basis (Hermite functions, Fourier modes on a box, wavelets).
The quantum theory \emph{counts}.
Its native cardinal is $\alef$, the infinity of the integers~\cite{cantor1891,halmos1974}.

General relativity (GR) is built on a smooth $4$-manifold $(M,g)$.
Any nonempty open $4$-manifold has $\abs{M}=\cont=2^{\alef}$ points.
A smooth metric is a section of $\mathrm{Sym}^{2}T^{*}M$; the space of such sections has cardinality $\cont$, because a continuous function is determined by its values on a countable dense set~\cite{cantor1878}.
The diffeomorphism group $\Diff(M)$ likewise has cardinality $\cont$.
The Einstein equation is a differential equation on this continuum.
GR's native cardinal is $\cont$, the infinity of the reals.

These are not aesthetic preferences.
They are the cardinalities of the objects the theories actually use.
The obstruction to a theory of everything is not only that one theory is quantum and the other is geometric.
It is that one theory lives at $\alef$ and the other lives at $\cont$, and Cantor's theorem says these are different sizes~\cite{cantor1891}:
\begin{equation}
  \alef \;<\; 2^{\alef} \;=\; \cont.
\end{equation}
There is no bijection between the integer skeleton of QFT and the continuum spacetime of GR.
Any unification must therefore \emph{supply a map} from one infinity to the other, and a map back.

The map forward is sampling plus inflation.
The map back is a weighted merge.
The weights are Kolmogorov complexity computed in a language that already contains Einstein's equation.
That is the whole theory.
The rest of this paper is the mathematics of those three sentences.

\begin{figure}[t]
\centering
{\small
\begin{tabular}{@{}c@{}}
$\ket{\Omega}\in\Hil$ \quad ($\alef$, QFT skeleton)\\[2pt]
$\big\downarrow$ \quad Born sample\\[2pt]
$s\in\{0,1\}^{*}$ \quad (integer Cauchy data)\\[2pt]
$\big\downarrow$ \quad GR inflation $\Gamma$\\[2pt]
$U=\Gamma(s)\in\Geo$ \quad ($\cont$, continuum geometry)\\[2pt]
$\big\downarrow$ \quad $\KG$\\[2pt]
$W(s)\propto \abs{\psi_{s}}^{2}\,2^{-\KG(s)}$\\[2pt]
$\big\downarrow$ \quad merge\\[2pt]
$\rho_{\star}=\sum_{s}W(s)\,\ket{U_{s}}\bra{U_{s}}$ \quad (countable sum)\\[2pt]
$\big\downarrow$ \quad effective state\\[2pt]
back to a vector in $\Hil$
\end{tabular}}
\caption{The $k$-cycle.
Quantum field theory supplies a countable state.
Sampling reads a finite string from that state.
General relativity completes the string to a continuum universe.
Gravitational Kolmogorov complexity weights the universe.
The merge is a countable sum of those weighted universes and returns an effective quantum state.}
\label{fig:cycle}
\end{figure}

\section{Cantor, Kolmogorov, Einstein}

\subsection{Cardinals}

Write $\abs{X}$ for the cardinality of a set $X$.
Cantor showed that $\abs{\N}=\alef$ is the smallest infinite cardinal, that $\abs{\R}=\cont=2^{\alef}$, and that there is no bijection $\N\to\R$~\cite{cantor1874,cantor1891}.
The identification $\cont=2^{\alef}$ is the statement that the continuum is (in bijection with) the power set of the integers: every real in $[0,1]$ is a subset of $\N$, via its binary expansion
\begin{equation}
  r \;=\; \sum_{n=1}^{\infty} b_{n}\,2^{-n},\qquad b_{n}\in\{0,1\}.
  \label{eq:binary}
\end{equation}
The continuum hypothesis $\mathrm{CH}$ asserts $\cont=\aleph_{1}$, i.e.\ that no cardinal sits strictly between $\alef$ and $\cont$.
$\mathrm{CH}$ is independent of ZFC~\cite{godel1940,cohen1963}.
Physics, in this theory, never uses an intermediate cardinal: samples are countable, completions are continuum, and nothing is sampled in between.
That is a physical reason to regard $\mathrm{CH}$ as true \emph{for physics}, whatever its status in set theory.
We return to this in Section~\ref{sec:ch}.

\begin{proposition}[Cardinalities of the two theories]
\label{prop:card}
Let $\Hil$ be an infinite-dimensional separable complex Hilbert space, and let $(M,g)$ be a smooth connected Lorentzian $4$-manifold.
Then:
\begin{enumerate}
  \item $\Hil$ admits a countable orthonormal basis, so the kinematic labels of QFT have cardinality $\alef$.
  \item As a set, $\abs{\Hil}=\cont$.
  \item $\abs{M}=\cont$, and the space $C^{\infty}(M,\mathrm{Sym}^{2}T^{*}M)$ of smooth metrics has cardinality $\cont$.
  \item The set of finite binary strings $\{0,1\}^{*}$ has cardinality $\alef$.
  \item The set of prefix-free programs for a fixed universal Turing machine has cardinality $\alef$.
\end{enumerate}
\end{proposition}

\begin{proof}
(1) is the definition of separability.
(2) follows because every vector is an $\ell^{2}$ sequence of complex amplitudes and $\abs{\C^{\N}}=\cont^{\alef}=(2^{\alef})^{\alef}=2^{\alef\cdot\alef}=2^{\alef}$.
(3): a nonempty open subset of $\R^{4}$ has cardinality $\cont$; a continuous real function on a second-countable manifold is determined by its restriction to a countable dense subset, so $\abs{C^{0}(M,\R)}\le\abs{\R^{\N}}=\cont^{\alef}=\cont$, and the smooth metrics inject into this set.
(4) and (5) are countable unions of finite sets.
\end{proof}

The distinction in (1) versus (2) is the entire subject.
QFT \emph{uses} a countable basis.
Its states, as points of $\Hil$, already have cardinality $\cont$---the same cardinal as GR geometries.
What QFT does \emph{not} have, and GR does, is a continuum of \emph{independent classical points}.
A quantum state is a square-summable sequence.
A spacetime is an uncountable set equipped with a smooth metric.
The map from the first to the second is not given by the Hilbert space structure alone.
It has to be put in by hand as a physical law.
That law is sampling plus Einstein inflation.

\subsection{Kolmogorov complexity}

Let $\mathsf{U}$ be a prefix-free universal Turing machine.
The prefix Kolmogorov complexity of a finite string $s\in\{0,1\}^{*}$ is~\cite{kolmogorov1965,chaitin1969,solomonoff1964,li2008}
\begin{equation}
  K(s) \;=\; \min\bigl\{\,\abs{p}\;:\; \mathsf{U}(p)=s\,\bigr\}.
\end{equation}
$K(s)$ is the length of the shortest program that prints $s$ and halts.
It is uncomputable, invariant up to an additive constant under change of $\mathsf{U}$ (the invariance theorem), and it induces the Solomonoff--Levin semimeasure~\cite{solomonoff1964,levin1974}
\begin{equation}
  \mathbf{M}(s) \;=\; \sum_{\mathsf{U}(p)=s} 2^{-\abs{p}} \;\ge\; 2^{-K(s)}.
  \label{eq:solomonoff}
\end{equation}
Kraft's inequality says $\sum_{p} 2^{-\abs{p}}\le 1$ for a prefix-free set of programs, so $\mathbf{M}$ is a semimeasure on strings: a universal Occam prior.
Schmidhuber proposed that a theory of everything \emph{is} a choice of such a prior on universe-histories~\cite{schmidhuber2000}.
Dzhunushaliev proposed that the probability of a gravitational configuration is $e^{-K[g]}$, and that the Einstein equation is the algorithm that reconstructs a $4$-geometry from initial data, thereby collapsing the complexity of Schwarzschild spacetime to the complexity of a single horizon radius plus the complexity of Einstein's equation~\cite{dzhunushaliev1998,dzhunushaliev2002}.

Those two observations are the seed of the weight in this theory.
What they do not yet supply is the cardinal split, the sampling map from a quantum state, or the merge.

\subsection{Einstein as a compressor}

The Cauchy problem of vacuum GR is locally well-posed~\cite{choquetbruhat1952,hawking1973}: a pair $(\gamma,K)$ of Riemannian $3$-metric and extrinsic curvature on a $3$-manifold $\Sigma$, satisfying the Hamiltonian and momentum constraints, determines a unique (up to diffeomorphism) maximal globally hyperbolic development $(M,g)$.
If the data are analytic, they are fixed by a countable list of Taylor coefficients on a dense set.
If the data are specified only to finite precision---and every sample is---they are a finite string.

\begin{remark}[GR is a compressor]
\label{rem:compressor}
A generic smooth metric on a compact region, written to precision $\varepsilon$ in Planck units, requires a number of bits that scales with the $4$-volume.
An Einstein solution on the same region is determined by $3$-data on a Cauchy slice, plus the Einstein equation.
The Einstein equation is a short program.
Therefore Einstein solutions are exponentially more compressible than generic metrics.
This is not an analogy with Kolmogorov complexity.
It \emph{is} Kolmogorov complexity, relative to a machine that may call an Einstein solver.
\end{remark}

Dzhunushaliev's estimate for Schwarzschild is the prototype~\cite{dzhunushaliev1998}: the whole $4$-geometry is recovered from $r_{g}/r_{\mathrm{Pl}}$ together with Einstein's equation, so
\begin{equation}
  K(\text{Schwarzschild}) \;\approx\; L\!\left[\Bigl(\frac{r_{g}}{r_{\mathrm{Pl}}}\Bigr)^{2}\right] + L_{\mathrm{Einstein}},
\end{equation}
where $L[\cdot]$ is program length for the indicated datum.
The Bekenstein--Hawking entropy $A/4\ell_{P}^{2}$ is the natural information content of horizon data~\cite{bekenstein1973,hawking1975}.
We will take this seriously as a theorem schema, not a coincidence (Section~\ref{sec:bh}).

\section{Postulates}

Five postulates.
Nothing else is put in.

\begin{postulate}[Integer quantum]
\label{post:qft}
The kinematical quantum state of the universe is a vector $\ket{\Omega}$ in a separable Hilbert space $\Hil$.
There exists a countable orthonormal basis $\{\ket{e_{n}}\}_{n\in\N}$ of $\Hil$, the \emph{integer skeleton}.
All quantum observables used in sampling are defined with respect to this skeleton or a uniformly equivalent one.
\end{postulate}

\begin{postulate}[Continuum gravity]
\label{post:gr}
A classical universe is a diffeomorphism class of globally hyperbolic Lorentzian $4$-manifolds with matter fields, $(M,g,\Phi)$.
Write $\Geo$ for the space of such classes.
$\Geo$ has cardinality $\cont$.
The dynamics on $\Geo$ is Einstein's equation coupled to the matter whose quantum version lives in $\Hil$.
\end{postulate}

\begin{postulate}[Sampling]
\label{post:sample}
A sample is a finite binary string $s\in\{0,1\}^{*}$ obtained by a finite-resolution measurement of $\ket{\Omega}$ on the integer skeleton.
The set of samples is countable.
Each sample $s$ is interpreted as Cauchy data (geometry plus matter) to finite precision.
The \emph{inflation map}
\begin{equation}
  \Gamma:\{0,1\}^{*} \;\longrightarrow\; \Geo
\end{equation}
sends $s$ to the maximal Einstein development of that data, or to a specified failure object if the constraints cannot be solved.
\end{postulate}

\begin{postulate}[Gravitational Kolmogorov weight]
\label{post:weight}
Let $\UTM$ be a prefix-free universal Turing machine whose language contains a primitive solver for the Einstein constraints and evolution (a \emph{GR-machine}).
The gravitational Kolmogorov complexity of a sample is
\begin{equation}
  \KG(s) \;=\; \min\bigl\{\,\abs{p}\;:\; \UTM(p)\ \text{outputs a description of}\ \Gamma(s)\ \text{to the resolution of}\ s\,\bigr\}.
\end{equation}
The unnormalized weight of $s$ is
\begin{equation}
  w(s) \;=\; \bigl|\braket{e_{s}}{\Omega}\bigr|^{2}\, 2^{-\KG(s)},
  \label{eq:weight}
\end{equation}
where $\ket{e_{s}}$ is the skeleton vector (or POVM element) that produced $s$.
\end{postulate}

\begin{postulate}[Merge]
\label{post:merge}
The physical universe is not a single sample.
It is the merge of all samples against the normalized weights
\begin{equation}
  W(s) \;=\; \frac{w(s)}{\sum_{s'} w(s')}.
  \label{eq:W}
\end{equation}
The merge has a coherent form (path integral), a decoherent form (ensemble), and a typical form (what an observer inhabits).
All three are defined by the same countable family $\{W(s),\Gamma(s)\}$.
\end{postulate}

The cycle is Figure~\ref{fig:cycle}.
QFT provides $\ket{\Omega}$ and the integers.
Sampling reads a string.
GR writes a continuum.
$\KG$ scores the writing.
The merge reads the scores and returns physics.

\section{\texorpdfstring{Sampling: from $\alef$ to $\cont$}{Sampling: from aleph\_0 to the continuum}}
\label{sec:sampling}

\subsection{What is sampled}

The vector $\ket{\Omega}$ is not itself a universe.
It is a countable list of amplitudes.
To obtain a universe one must extract a definite string.

Let $\{E_{s}\}_{s\in\{0,1\}^{\le N}}$ be a finite-resolution POVM on $\Hil$ indexed by strings of length at most $N$, refining as $N\to\infty$ to a complete measurement in a skeleton basis.
A sample of resolution $N$ is a draw $s$ with Born probability
\begin{equation}
  \mathrm{Pr}(s\mid\Omega,N) \;=\; \tr\bigl(E_{s}\,\ket{\Omega}\bra{\Omega}\bigr).
\end{equation}
In a rank-one skeleton measurement this is $\abs{\braket{e_{s}}{\Omega}}^{2}$.
The resolution $N$ is Planck-scale information capacity, not an observer's laboratory precision: it is the number of bits the integer skeleton offers before inflation.

The string $s$ is Cauchy data.
Concretely, one may take $s$ to encode:
\begin{itemize}
  \item a finite simplicial $3$-complex, or a finite causal set, or a finite graph with holonomies (a spin network), together with
  \item matter field values on the vertices, to a finite number of bits, and
  \item a finite specification of lapse and shift, or a statement that they are to be solved for.
\end{itemize}
The theory does not prefer a discretization.
Any countable encoding of $3$-data is a choice of skeleton.
Uniformly equivalent encodings change $K$ by $O(1)$ and leave the merge invariant at leading exponential order.

\subsection{Inflation is Cantor's map}

The inflation map $\Gamma$ is the physical form of~\eqref{eq:binary}.
A binary sequence is a real.
A countable list of Cauchy data is a spacetime, once Einstein's equation is allowed to run.

\begin{definition}[Inflation]
If $s$ encodes constraint-satisfying data $(\gamma,K,\Phi)$ on a $3$-manifold, $\Gamma(s)$ is a representative of the diffeomorphism class of the maximal globally hyperbolic development.
If the constraints fail, $\Gamma(s)$ is the development of the nearest constraint-satisfying projection of the data (in a fixed Sobolev metric), and the failure is charged to $\KG$ as extra program length; see~\eqref{eq:KG-split}.
If the development is incomplete in a non-physical way (naked singularity, closed timelike curves), $\Gamma$ returns that incomplete object and $\KG$ pays for the pathology.
\end{definition}

Two structural facts make this a map $\alef\to\cont$ rather than a trick.

First: the domain is countable.
There are $\alef$ finite strings.
There are $\alef$ programs.
Sampling cannot emit an uncountable object.
Whatever continuum structure a universe has, it must be \emph{generated} from a countable seed.

Second: the Einstein Cauchy problem is a completion.
It takes data on a set that can be chosen countable-dense (analytic data; or a finite net, in the finite-resolution case) and produces a metric at every point of a $4$-manifold of cardinality $\cont$.
This is the same logical shape as completing $\Q$ to $\R$, or taking a power set.
The continuum of GR is not an independent infinity floating beside QFT.
It is the Cantor image of QFT's integers, under the map ``solve Einstein.''

\begin{proposition}[Cantor inflation]
\label{prop:cantor-inf}
The image $\Gamma(\{0,1\}^{*})$ is a countable subset of $\Geo$.
Its metric completion, in any topology in which Einstein developments of refining data converge, has cardinality at most $\cont$ and contains the classical spacetimes that finite-resolution observers can stably identify.
\end{proposition}

The physical continuum is a completion of a countable set of inflations, just as $\R$ is a completion of $\Q$.
Generic elements of $\Geo$---metrics that are not Einstein developments of any finite string---have $\KG=\infty$ in the sense that no program of the GR-machine prints them to arbitrary resolution.
They have Solomonoff weight zero.
They are the non-measurable, non-computable geometries that the gravitational path integral should never have been summing over.

\subsection{The quantum state of the universe}

Postulate~\ref{post:qft} leaves $\ket{\Omega}$ unspecified up to the demand that it live in $\Hil$.
A theory of everything must pick one.

\begin{definition}[Kolmogorov vacuum]
\label{def:omega}
$\ket{\Omega}$ is a diffeomorphism-invariant state of least gravitational Kolmogorov complexity among states that admit a Fock representation of the matter fields on some Einstein background in the support of $W$.
Equivalently, $\ket{\Omega}$ is a shortest program that prints a Hartle--Hawking no-boundary state~\cite{hartle1983}: compact Euclidean $4$-geometries with a single boundary, no other boundary.
\end{definition}

The no-boundary proposal is already a short program.
That is why it is a candidate.
The tunneling proposal~\cite{vilenkin1983} is another short program, of comparable length.
The merge does not require us to decide between them at the level of postulates: both are samples, both receive weights, and $W$ concentrates on whichever is shorter once the Einstein solver is in the language.
If they have comparable $\KG$, both survive in the decoherent merge as alternative cosmological branches.

This is the sense in which the theory does not put in a wave function of the universe by hand.
It puts in a universal machine with Einstein's equation, and $\ket{\Omega}$ is a shortest vector that machine can name.

\section{\texorpdfstring{Weight: GR/$k$-complexity}{Weight: GR/k-complexity}}
\label{sec:weight}

The user of a theory of everything is entitled to one number per universe.
That number is $w(s)$ in~\eqref{eq:weight}.
It has two factors, and they are not independent.

\subsection{The Born factor}

$\abs{\braket{e_{s}}{\Omega}}^{2}$ is the quantum sampling measure.
It is the probability that the integer skeleton, interrogated at finite resolution, returns $s$.
Without this factor the theory would be a classical algorithmic cosmology in the style of Schmidhuber~\cite{schmidhuber2000}: possible, but silent about interference, identical particles, and the observed linearity of quantum amplitudes.
With this factor, sampling is a measurement of $\ket{\Omega}$.
The many worlds of Everett~\cite{everett1957} are the samples.
The present theory adds a second weight on those worlds, and a merge.

\subsection{The gravitational Kolmogorov factor}

$2^{-\KG(s)}$ is Occam's razor computed on a GR-machine.
Three comments.

\paragraph{Why a GR-machine, not a bare Turing machine.}
On a bare machine, a spacetime is a huge array of metric components, and $K(s)$ is essentially the raw bit-length of that array unless a compression happens to exist.
Einstein's equation \emph{is} that compression, but only if the machine is allowed to use it.
Putting the Einstein solver in the language of $\UTM$ is the statement that GR is fundamental, not emergent from a GR-blind program search.
(If GR were itself a short program on a bare machine, the invariance theorem would absorb it into an $O(1)$ and the two choices would agree.
If it is not, the GR-machine is a different theory, and it is the theory we mean.)

\paragraph{The split of $\KG$.}
Write, schematically,
\begin{equation}
  \KG(s) \;=\; K(s_{\mathrm{data}}\mid \text{Einstein}) \;+\; \frac{I_{\mathrm{off}}[\Gamma(s)]}{\hbar\ln 2} \;+\; O(1).
  \label{eq:KG-split}
\end{equation}
The first term is the program length of the Cauchy data given that Einstein evolution is free.
The second term charges constraint violation and off-shellness: $I_{\mathrm{off}}$ is a non-negative functional that vanishes on Einstein solutions and grows with $\int\abs{G_{\mu\nu}-8\pi T_{\mu\nu}}^{2}$.
On shell, $I_{\mathrm{off}}=0$ and $\KG$ reduces to the complexity of the data plus the constant $L_{\mathrm{Einstein}}$.
Off shell, the extra bits kill the weight exponentially.

\paragraph{Action as complexity.}
The Euclidean Einstein--Hilbert action
\begin{equation}
  I_{E}[g] \;=\; -\frac{1}{16\pi G}\int_{M}\!R\,\sqrt{g}\,\dd^{4}x \;-\; \frac{1}{8\pi G}\int_{\partial M}\!K\,\sqrt{h}\,\dd^{3}x
\end{equation}
is the continuum shadow of $\KG$.
Hartle and Hawking assign a universe the amplitude $e^{-I_{E}/\hbar}$~\cite{hartle1983}.
Dzhunushaliev replaces the action in the path integral by $K[g]$~\cite{dzhunushaliev1998}.
The present theory says these are the same assignment once units are fixed by
\begin{equation}
  \frac{I_{E}[\Gamma(s)]}{\hbar} \;\sim\; \KG(s)\,\ln 2 \qquad \text{on the dominant saddles}.
  \label{eq:action-K}
\end{equation}
This is a matching condition, not an identity for every metric.
It is the statement that the semiclassical gravitational measure and the Solomonoff measure agree where both are reliable: on Einsteinian saddles of low complexity.
Away from those saddles Solomonoff is the definition, and the continuum action is an approximation.

\begin{proposition}[Einsteinian dominance]
\label{prop:einstein-dom}
Let $s_{\mathrm{E}}$ encode constraint-satisfying data of bit-length $n$, and let $s_{\mathrm{rand}}$ encode a generic metric array of the same resolution, not an Einstein development of any shorter string.
Then
\begin{equation}
  \KG(s_{\mathrm{E}}) \;\le\; n + L_{\mathrm{Einstein}} + O(1), \qquad \KG(s_{\mathrm{rand}}) \;\ge\; n - O(\log n),
\end{equation}
and the weight ratio satisfies
\begin{equation}
  \frac{w(s_{\mathrm{rand}})}{w(s_{\mathrm{E}})} \;\le\; 2^{-L_{\mathrm{Einstein}}+O(\log n)}\cdot \frac{\abs{\braket{e_{s_{\mathrm{rand}}}}{\Omega}}^{2}}{\abs{\braket{e_{s_{\mathrm{E}}}}{\Omega}}^{2}}.
\end{equation}
If $\ket{\Omega}$ is not conspiratorially supported on random geometries, Einsteinian samples dominate the merge by a factor exponential in the program length of Einstein's equation.
\end{proposition}

That is why classical GR exists.
It is not because the continuum was there first.
It is because the shortest programs that print a $4$-geometry print Einstein solutions.

\subsection{Normalization}

The denominator of~\eqref{eq:W} is a countable sum of non-negative terms.
It may diverge.
If it does, one passes to a sequence of finite-resolution cutoffs $N$ (strings of length $\le N$) and takes Banach limits of the finite-$N$ weights, or equivalently uses the Solomonoff semimeasure without forcing $\sum W=1$.
This is standard~\cite{li2008}.
The gravitational path integral has always had a measure problem~\cite{gibbons1979}.
Here the domain of summation is $\alef$, not $\cont$, so the only remaining issue is the possible divergence of a series of positive terms, which is an ordinary analytic problem, not a foundational one.

\section{\texorpdfstring{Merge: from $\cont$ back to $\alef$}{Merge: from the continuum back to aleph\_0}}
\label{sec:merge}

Sampling produces a cloud of continuum universes, each tagged with a weight $W(s)$.
Physics is what one does with the cloud.
There are three consistent operations, and they are the three regimes of the theory.

\subsection{Coherent merge: the $k$-integral}

The coherent merge is the replacement for the gravitational path integral:
\begin{equation}
  Z[J] \;=\; \sum_{s\in\{0,1\}^{*}} \bigl|\braket{e_{s}}{\Omega}\bigr|^{2}\, 2^{-\KG(s)}\, \exp\Bigl(\frac{i}{\hbar}S_{\mathrm{matter}}[\Gamma(s);J]\Bigr).
  \label{eq:Z}
\end{equation}
It is a countable sum.
There is no continuum measure $\mathcal{D}g$ over metrics, no Faddeev--Popov determinant waiting on a non-existent Lebesgue measure on $\Geo$, no conformal-factor problem of the Euclidean integral as a \emph{definition} (it may reappear as a property of the dominant samples).
The continuum is inside $\Gamma(s)$, not in the domain of summation.

When $\KG$ is matched to $I_{E}$ as in~\eqref{eq:action-K} and the Born factor is slowly varying,~\eqref{eq:Z} reduces to a skeleton-regularized Hartle--Hawking integral.
The present definition is more general: it assigns weight to off-shell samples through $\KG$, and it knows about the quantum state $\ket{\Omega}$.

The effective amplitude for a coarse-grained $3$-geometry $h$ is
\begin{equation}
  \Psi[h] \;=\; \sum_{s:\,\Gamma(s)|_{\Sigma}\simeq h} \sqrt{w(s)}\, e^{iS[\Gamma(s)]/\hbar}.
  \label{eq:Psi}
\end{equation}
This $\Psi$ is the wave function of the universe as an output, not an input.
It is annihilated by the constraints to the extent that $\UTM$ treats diffeomorphisms as free (they do not add to $\KG$), which is the path-integral derivation of Wheeler--DeWitt~\cite{halliwell1991} rewritten in Kolmogorov language.

\subsection{Decoherent merge: the ensemble}

Samples that differ by a macroscopic geometry do not interfere.
The decoherence functional of Gell-Mann and Hartle~\cite{gellmann1993} is, in this theory, a statement about $\KG$:
\begin{equation}
  D(s,s') \;=\; \sqrt{W(s)W(s')}\, \exp\bigl(-K(s:s')\bigr),
\end{equation}
where $K(s:s')$ is the mutual algorithmic information.
When two inflations $\Gamma(s)$ and $\Gamma(s')$ have macroscopic metric differences, a short program cannot convert one into the other, $K(s:s')$ is large, and $D$ is diagonal.
The decoherent merge is then the mixed state
\begin{equation}
  \rho_{\star} \;=\; \sum_{s} W(s)\, \ket{U_{s}}\bra{U_{s}}, \qquad \ket{U_{s}}\;=\;\text{a code of}\ \Gamma(s).
  \label{eq:rho}
\end{equation}
Classical spacetime is a diagonal ensemble of Einsteinian inflations.
This is also the sense in which the theory is a many-worlds theory with a preferred measure: the worlds are the samples, the measure is $W$, and the preferred basis is the one that diagonalizes $D$, i.e.\ the low-mutual-information inflations.

\subsection{Typical merge: the inhabited universe}

An observer does not see $\rho_{\star}$.
An observer inhabits a single inflation $\Gamma(s_{\star})$, where $s_{\star}$ is $W$-typical:
\begin{equation}
  s_{\star} \;\in\; \bigl\{\,s\;:\; -\log_{2} W(s) \;\le\; \KG(s) + K(\Omega) + O(1)\,\bigr\}.
\end{equation}
Typicality in the Solomonoff measure is the precise form of ``the simplest universe consistent with the quantum state.''
The inhabited geometry $\bar{g}=\Gamma(s_{\star})$ is the background of effective field theory.
QFT, in the ordinary textbook sense, is the residual integer dynamics of $\ket{\Omega}$ in the Fock representation on $(\bar{M},\bar{g})$.

Operationally, every prediction is a weighted expectation
\begin{equation}
  \langle\mathcal{O}\rangle \;=\; \sum_{s} W(s)\, \mathcal{O}\bigl[\Gamma(s)\bigr].
  \label{eq:expect}
\end{equation}
When $W$ concentrates on a tight family of diffeomorphic Einstein solutions,~\eqref{eq:expect} is the classical value of $\mathcal{O}$ on that solution, plus QFT fluctuations from the residual state.
When $W$ is spread (near a topology change, near a black-hole singularity, in the Planck regime),~\eqref{eq:expect} is a genuine ensemble average and classical geometry fails.

\subsection{The merge returns a quantum state}

The output of the merge is again an object in a separable Hilbert space: $\rho_{\star}$ in~\eqref{eq:rho}, or $\Psi$ in~\eqref{eq:Psi}.
The cycle in Figure~\ref{fig:cycle} closes.
This is required for consistency.
If the merge produced a naked continuum object with no countable description, the theory would have used $\cont$ as an output while forbidding it as an input.
Instead:
\begin{equation}
  \alef \;\xrightarrow{\ \Gamma\ }\; \cont \;\xrightarrow{\ W\ }\; \alef.
\end{equation}
The continuum is an intermediate, generated and then integrated out.
That is what a well-defined gravitational path integral always should have been.

\section{Recovery of general relativity}
\label{sec:gr}

Classical GR is the typical merge of low-$\KG$ samples.

\begin{enumerate}
  \item \textbf{Einstein's equation.}
  Proposition~\ref{prop:einstein-dom}: the GR-machine prints Einstein developments far more cheaply than generic metrics, so $W$ concentrates on $\mathrm{Ein}(\Geo)\subset\Geo$.
  Off-shell samples survive only at Planckian resolution, where $I_{\mathrm{off}}/\hbar$ is not large.

  \item \textbf{Diffeomorphism invariance.}
  A diffeomorphism is a change of coordinates.
  On the GR-machine it is a free subroutine: applying a diffeomorphism to a named metric does not increase $\KG$ by more than $O(1)$ plus the complexity of the diffeomorphism itself, and the diffeomorphisms that are namable by short programs are the ones physics already treats as gauge.
  Two samples that inflate to diffeomorphic spacetimes are identified in $\Geo$ (Postulate~\ref{post:gr}), so they are the same universe.

  \item \textbf{Local well-posedness and causality.}
  Inflation \emph{is} the Cauchy problem.
  Globally hyperbolic developments are exactly the objects $\Gamma$ produces on shell.
  Closed timelike curves require extra data (identifications) that increase $\KG$; they are not forbidden by fiat, they are expensive.

  \item \textbf{The Einstein--Hilbert action as Hamilton--Jacobi function.}
  On the dominant samples,~\eqref{eq:action-K} holds, so the phase of the coherent merge is the Einstein--Hilbert phase, and the WKB limit of $\Psi[h]$ is classical GR with that action.
  Lovelock uniqueness~\cite{lovelock1971} says that in four dimensions the only second-order diffeomorphism-invariant tensor built from the metric is Einstein's.
  A GR-machine whose Einstein solver was replaced by a higher-order PDE would have a longer program (more jets to specify) or would lose well-posedness (and then pay $I_{\mathrm{off}}$).
  Four-dimensional Einstein gravity is the shortest well-posed compressor of $3$-data into $4$-geometry.

  \item \textbf{Matter coupling.}
  The matter content of $\Gamma(s)$ is whatever field data $s$ carried, evolved by the coupled Einstein--matter system.
  The shortest matter programs compatible with the residual quantum state are the ones that appear in the typical universe.
  This is not a derivation of the Standard Model.
  It is a statement of what would count as a derivation: a demonstration that the Standard Model Lagrangian is a shortest matter plugin for $\UTM$ consistent with $\ket{\Omega}$ and with the observed low-energy sample.
\end{enumerate}

\section{Recovery of quantum field theory}
\label{sec:qft}

QFT is the integer theory on a merged background.

Fix a typical inhabited geometry $(\bar{M},\bar{g})=\Gamma(s_{\star})$.
The residual quantum state is the condition of $\ket{\Omega}$ on the skeleton, given $s_{\star}$.
In a neighbourhood of $s_{\star}$ in which inflation produces geometries in which a Fock representation exists (Hadamard states on globally hyperbolic backgrounds~\cite{wald1994,kay1991}), this residual state \emph{is} a QFT state of the matter fields, and of gravitons as linearized excitations of $\bar{g}$.

\begin{enumerate}
  \item \textbf{Particles are integers.}
  Occupation numbers, species indices, momenta in a box: $\alef$.
  There is no continuum of particle types because a continuum of types would be a $\cont$-sized label set, which sampling cannot emit.
  A ``continuous spin'' or a continuous family of species can appear only as a classical limit of many discrete samples.

  \item \textbf{Fields are continuum completions.}
  A quantum field operator is indexed by test functions on $\bar{M}$, and $\bar{M}$ has cardinality $\cont$.
  This is inflation again: the field as a distribution is the GR-plus-matter completion of countable mode data.
  The notorious infinities of QFT are $\alef$-sums (mode sums, loop momenta on a lattice going to infinity) pretending to be $\cont$-integrals without having passed through the merge.
  The merge caps them.

  \item \textbf{Ultraviolet divergences and holography.}
  A local QFT Hilbert space is infinite-dimensional because a region of continuum spacetime supports $\alef$ modes.
  Gravity does not permit this.
  The Bekenstein bound and its covariant forms~\cite{bekenstein1981,bousso1999}, and the independent argument that a local region in quantum gravity has finite-dimensional Hilbert space~\cite{bao2017}, say that a region of boundary area $A$ supports at most $A/4\ell_{P}^{2}$ independent bits.
  In the present theory that statement is the assertion
  \begin{equation}
    \KG(s|_{R}) \;\le\; \frac{A(\partial R)}{4\ell_{P}^{2}} + O(\log A).
    \label{eq:holo}
  \end{equation}
  A sample that tries to specify more independent data in $R$ than the area allows either inflates to a geometry in which $R$ has already collapsed to a black hole (and the extra data are behind a horizon, charged to a different region), or it is not a shortest program and is exponentially downweighted.
  The continuum interior of $R$ is a $\Gamma$-completion of a countable boundary code of length $\sim A/4\ell_{P}^{2}$.
  That is holography~\cite{tHooft1993,susskind1995,maldacena1998} as Cantor inflation with a Kraft bound.

  \item \textbf{The Born rule.}
  Laboratory Born frequencies are the $W$-weights of samples that differ only in a small subsystem, with the environment and the geometry held fixed.
  For those samples $\KG$ is essentially constant (the extra bits specifying which detector clicked are $O(1)$ compared with the complexity of the apparatus), so $W(s)\propto\abs{\psi_{s}}^{2}$.
  Corrections are considered in Section~\ref{sec:predictions}.

  \item \textbf{QFT in curved spacetime.}
  The typical background $\bar{g}$ need not be Minkowski.
  The residual state on $\bar{g}$ is ordinary QFT in curved spacetime~\cite{birrell1982,wald1994}.
  Hawking radiation is the residual integer dynamics on a collapsing typical sample~\cite{hawking1975}.
  Inflationary perturbations are the residual dynamics on a typical inflating sample~\cite{mukhanov1992}.
\end{enumerate}

\section{Black holes}
\label{sec:bh}

A Schwarzschild sample of mass $M$ is a very short string: one real number $r_{g}=2GM$, plus Einstein's equation, plus a statement of vacuum and spherical symmetry~\cite{dzhunushaliev1998}.
Spherical symmetry is itself a short program (``all points at radius $r$ are equivalent'').
Hence Schwarzschild, Kerr, and FLRW are not generic points of $\Geo$.
They are the Kolmogorov elite.
That is why they are the first solutions anyone writes down, and why they dominate the merge among macroscopic isolated objects.

The horizon is the place where the integer skeleton of the exterior ends.
A sample of the exterior of a black hole of area $A$ cannot cheaply specify independent interior data beyond $A/4\ell_{P}^{2}$ bits:~\eqref{eq:holo} applied to the exterior region.
The gravitational Kolmogorov complexity of the hole, as an object in the merge, is therefore
\begin{equation}
  \KG(\text{BH of area }A) \;=\; \frac{A}{4\ell_{P}^{2}} + L_{\mathrm{Einstein}} + K(\text{charges, angular momenta}) + O(\log A).
  \label{eq:KBH}
\end{equation}
The first term is Bekenstein--Hawking entropy~\cite{bekenstein1973,hawking1975,gibbons1977}.
In this theory it is not a thermodynamic analogy.
It is the program length of the shortest GR-machine description of the hole.

The decoherent merge of samples that agree on the exterior and differ on the interior by up to $e^{A/4\ell_{P}^{2}}$ bits is a mixed state of von Neumann entropy $A/4\ell_{P}^{2}$.
That mixed state is the Hawking radiation thermal state, viewed from the outside.
Information is not lost in the coherent merge: the coherent sum~\eqref{eq:Z} still contains the interior samples.
It is lost, operationally, in the decoherent merge, because $K(s_{\mathrm{ext}}:s_{\mathrm{int}})$ is large and the off-diagonal terms in $D$ vanish.
Unitary evaporation is the statement that at late times a new family of samples---the radiation---becomes a shortest description of the same universe, with $\KG$ matching~\eqref{eq:KBH}.
Whether that matching is exact (no firewall, no remnant) is a question about the GR-machine's shortest programs for evaporating geometries, not a question that has to be settled before the theory can be stated.

\section{Cosmology and the arrow of time}
\label{sec:cosmo}

\subsection{The typical cosmos}

A typical sample under $W$ has:
\begin{itemize}
  \item short Cauchy data (low $K(s_{\mathrm{data}})$),
  \item on-shell Einstein inflation ($I_{\mathrm{off}}=0$),
  \item a residual quantum state close to the Kolmogorov vacuum.
\end{itemize}
Short, constraint-satisfying, homogeneous data on a $3$-sphere or a torus, plus a short inflaton potential, are extremely cheap programs that inflate to large, smooth, approximately FLRW universes.
This is why the observed cosmos is large, smooth, and approximately isotropic.
It is not because a continuum of initial conditions was fine-tuned.
It is because the alternatives are longer programs.

Inflation is preferred as a compressor: a few e-folds of a short potential convert Planckian data into a vast Einsteinian region, which would otherwise have to be specified bit by bit.
Too much inflation, of a kind that requires a long fine-tuned potential or an enormous number of named e-folds, costs $\KG$.
The measure $W$ therefore favours a moderate inflating family, in the same qualitative direction as the Hartle--Hawking measure~\cite{hartle2008} and opposite to volume-weighting that rewards $10^{10^{10}}$ e-folds.
A quantitative $N_{e}$ is a computation of $\KG$ for concrete inflaton plugins; it is not settled by the postulates.

\subsection{The cosmological constant}

A generic vacuum energy is an incompressible real: specifying $\Lambda$ to $N$ bits costs $N$ bits.
A vacuum energy that is a short program---$0$, or a topological invariant, or a ratio of two short geometric quantities---is cheap.
The typical merge therefore does not produce a Planck-scale cosmological constant as a random QFT zero-point sum.
That sum is an $\alef$-divergence that the merge never performs (Section~\ref{sec:qft}).
What it produces is a $\Lambda$ with a short description, possibly zero, possibly the residual of a short compactification, possibly the smallest nonzero value namable by the matter plugin.
The observed value $\Lambda\sim 10^{-122}$ in Planck units is a $122$-bit number and is \emph{not} short as a raw integer.
The theory's claim is that the shortest program for the observed $\Lambda$ is not ``print this $122$-bit integer'' but a short geometric or vacuum-structure program whose output happens to be that integer.
If no such program exists, the observed $\Lambda$ is a genuine cost, and the typicality of our sample has to be paid for elsewhere (anthropic conditioning on the residual state, which the theory permits but does not advertise).

\subsection{The arrow of time}

Kolmogorov complexity of a typical inhabited universe, as a function of Cauchy time, is low on the initial slice---that is what ``typical under $W$'' means---and high later, because Einstein evolution of even short data produces geometrically rich $4$-regions whose uncompressed description is long.
The algorithmic entropy of the slice~\cite{zurek1989} increases.
That is the arrow of time.
It does not require a special postulate about past hypotheses beyond the statement that $W$ itself is the past hypothesis: the measure is supported on short data.

Boltzmann brains are long incompressible strings that fluctuate into existence in a typical late-time geometry.
Their $\KG$ is of order their own bit-length, so $W$ suppresses them exponentially relative to ordinary cosmological observers, who are cheap by-products of a short inflating program~\cite{carroll2017}.
This is the standard algorithmic resolution, now sitting on the same weight that defines the theory.

\section{The continuum hypothesis, physically}
\label{sec:ch}

Sampling emits elements of a countable set.
Inflation produces objects of cardinality $\cont$.
The merge sums over a countable index.
No construction in the theory produces a set of cardinality $\kappa$ with $\alef<\kappa<\cont$.

\begin{proposition}[Physical CH]
\label{prop:ch}
Every set that the $k$-cycle constructs from $\ket{\Omega}$ by sampling, inflation, weighting, and merging is either finite, of cardinality $\alef$, or of cardinality $\cont$ (or, if one takes power sets of geometries, of cardinality $2^{\cont}$, which the cycle never samples).
In particular the theory has no physical use for a cardinal strictly between $\alef$ and $\cont$.
\end{proposition}

This is not a proof of $\mathrm{CH}$ in ZFC, which is impossible~\cite{godel1940,cohen1963}.
It is a statement that physics, if it is this theory, is a model of set theory in which the only infinite cardinals that occur as types of physical objects are $\alef$ and $\cont$.
Intermediate cardinals, if they exist, are empirically idle.
That is the strongest physical meaning $\mathrm{CH}$ can have.

It is also a diagnosis of a class of failed unifications.
A theory that quantizes GR by putting a Hilbert space of dimension $\cont$ (a non-separable kinematical space of ``all metrics'') on top of the continuum, without a sampling map back to $\alef$, has used a cardinal that QFT does not know how to measure.
Loop quantum gravity's kinematical Hilbert space is non-separable in the standard construction and becomes separable once diffeomorphism moduli are treated as unphysical~\cite{fairbairn2004}; that move is, in the present language, a refusal to treat a $\cont$-sized label set as quantum kinematics.
The refusal is correct.

\section{What the theory is not}

It is not a derivation of the Standard Model spectrum, gauge group, or three generations.
Those are matter plugins to $\UTM$.
The theory says how a plugin is weighted, not which plugin is shortest.

It is not a claim that Kolmogorov complexity is computable, or that observers compute $\KG$ in order to exist.
$W$ is a defining measure, as the Euclidean action is a defining measure in Hartle--Hawking.
Predictions that require the exact value of $K$ are not available; predictions that require only crude bounds (Einsteinian dominance, holographic caps, shortness of Schwarzschild, suppression of Boltzmann brains) are.

It is not hypercomputation.
Section~\ref{sec:hyper} is the argument: classical GR admits Malament--Hogarth worldlines that decide the halting problem, hence compute $K$, but those geometries are long programs, are not globally hyperbolic, and are suppressed by $W$.
The GR-machine remains a Turing machine with an Einstein primitive.

It is not, despite the countable sums, a causal-set theory~\cite{bombelli1987}, a CDT theory~\cite{ambjorn2005}, or a Wolfram model~\cite{wolfram2020}, although each of those is a legitimate choice of integer skeleton.
The ontology is the pair $(\Hil,\Geo)$ with maps $\Gamma$ and $W$, not a particular discretization.

It is not a solution of Gödel--Chaitin incompleteness~\cite{chaitin1974,faizal2025}.
$\KG$ is uncomputable, $\ket{\Omega}$ is named but not constructed, and there will be physical questions whose answers are not theorems of any recursive axiomatization of the GR-machine.
The theory is complete as an ontology and a measure.
It is incomplete as a decision procedure.
That is the correct amount of completeness for a theory that takes Kolmogorov complexity as fundamental.

\section{Hypercomputation}
\label{sec:hyper}

Kolmogorov complexity is uncomputable, but only just.
The predicate $\KG(s)=n$ is $\Delta^{0}_{2}$: there exists a program of length $n$ that prints $s$, and none shorter.
The first half is the halting problem; the second is its complement.
So $K$ is computable from a halt oracle (Turing degree $0'$)~\cite{li2008}.
Any device that decides HALT computes $K$, by dovetailing all programs of length at most $n$ and asking which halt with output $s$.

Classical general relativity supplies such a device, on paper.

\subsection{Malament--Hogarth geometries}

A Malament--Hogarth spacetime contains a worldline $\lambda$ of infinite proper time that lies entirely in the causal past of an event $p$ reached by an observer in finite proper time~\cite{earman1993,hogarth2004}.
Send a Turing machine along $\lambda$.
If it ever halts, it signals to $p$.
The observer, arriving at $p$, either has the signal or does not.
That decides HALT, hence decides $K$.

Hogarth showed that a suitable MH geometry can settle any arithmetic sentence~\cite{hogarth2004}.
Etesi and N\'emeti argued that Kerr interiors implement a fragment of this~\cite{etesi2002,nemeti2006}.
Welch showed that MH geometries can resolve hyperarithmetic queries, with an ordinal bound $w(M)$ that is a constant of the manifold~\cite{welch2008}.
In the two-infinities language of this paper: an infinite integer computation is packed onto a continuum worldline and read off at one event.
That is $\alef$ being decided by $\cont$.
Inflation $\Gamma$ is the same logical shape---countable Cauchy data completed to a continuum spacetime.
Closed timelike curves are a second, different GR loophole; they are not needed for the MH construction.

\subsection{Why typical $k$-universes do not compute $W$}

Four independent blocks sit between that construction and ``observers evaluate $\KG$.''

\paragraph{The GR-machine is not an MH computer.}
Postulate~\ref{post:weight} equips $\UTM$ with an Einstein primitive, not with an infinite-time observer.
Finite-resolution numerical GR is Turing-computable, so $\KG$ differs from ordinary $K$ by $O(1)$ (the invariance theorem).
Inflation $\Gamma$ \emph{names} an Einstein development; it does not run a supertask along $\lambda$.

\paragraph{MH geometries are not typical under $W$.}
They fail global hyperbolicity~\cite{hogarth2004,earman1993}.
Inflation in this theory prefers globally hyperbolic Cauchy developments (Postulate~\ref{post:sample})~\cite{choquetbruhat1952}.
An MH geometry is a long program: extra identifications, an inner horizon or naked singularity, a computer worldline of infinite proper time.
That cost sits in $\KG$.
Einsteinian dominance (Proposition~\ref{prop:einstein-dom}) then suppresses it by $\sim 2^{-c}$ relative to a short globally hyperbolic sample.
The universes in which GR hypercomputes $K$ are exactly the universes the merge downweights.

\paragraph{Holography forbids the infinite tape.}
A region of area $A$ supports at most $A/4\ell_{P}^{2}$ independent bits~\cite{bekenstein1981,bousso1999,bao2017}.
An MH computer needs unbounded memory along $\lambda$.
Either those bits live in the observer's causal diamond, violating~\eqref{eq:holo}, or they live behind a horizon and cannot be communicated as a clean halt bit.
Infinite blueshift at the MH event~\cite{earman1993,etesi2002} is the same obstruction in GR language: the signal that would carry the extra information destroys the receiver.

\paragraph{An oracle does not give $\KG$ of the inhabited sample.}
Chaitin: a system of complexity $n$ cannot prove $K(s)>n+O(1)$ for strings it can name~\cite{chaitin1974}.
An observer whose own description is part of the sample cannot use an MH experiment inside that sample to certify $\KG$ of the whole sample.
They can at best compute $K$ of \emph{smaller} strings, the way a halt oracle computes $K$ of programs shorter than the oracle's own specification.
$W$ is a defining measure, like the Euclidean action in Hartle--Hawking.
It is not a quantity the universe has to calculate in order to exist.

\begin{proposition}[No typical halt oracle]
\label{prop:no-halt}
In a $W$-typical sample, there is no operational halt oracle available to a finite-area observer.
Classical GR's Malament--Hogarth constructions, where they exist as solutions, have Solomonoff weight exponentially small in their extra program length relative to a globally hyperbolic Einsteinian development of comparable Cauchy data.
\end{proposition}

The continuum of GR can look like a halt oracle.
The merge does not sample that oracle.
If a $W$-typical universe presented an observer with a usable Malament--Hogarth worldline---a clean halt bit from an infinite computation in finite proper time, without destroying the receiver---that would falsify Einsteinian dominance or the holographic bound, and so falsify the theory.
Hypercomputation in a typical sample would kill $k$, not complete it.

\section{Predictions and tests}
\label{sec:predictions}

A theory of everything earns the name by recovering the two existing theories and then saying something they do not.
The recoveries are Sections~\ref{sec:gr}--\ref{sec:cosmo}.
The extras are these.

\paragraph{1. The gravitational path integral is a countable sum.}
Any continuum measure $\mathcal{D}g$ is an approximation to~\eqref{eq:Z}, valid on Einsteinian saddles and nowhere a definition.
Practical consequence: lattice, spin-foam, and causal-dynamical-triangulation sums are closer to the definition than continuum Euclidean integrals are.
Their continuum limits, when they exist, are inflations, not domains of integration.

\paragraph{2. A complexity correction to Born frequencies.}
For two laboratory outcomes $a,b$ of a measurement on a small system,
\begin{equation}
  \frac{W(a)}{W(b)} \;=\; \frac{\abs{\psi_{a}}^{2}}{\abs{\psi_{b}}^{2}}\, 2^{-\bigl(\KG(a)-\KG(b)\bigr)}.
  \label{eq:born-corr}
\end{equation}
In a laboratory whose environment has already contributed $\sim 10^{23}$ bits to $\KG$, the difference $\KG(a)-\KG(b)$ is $O(1)$ and the correction is unobservable.
In a fully specified, isolated, few-qubit device whose total $\KG$ is dominated by the register, a systematic bias toward the lower-complexity outcome is in principle present.
This is the most direct empirical signature the theory has, and the easiest to overclaim.
It should be looked for only in experiments where the Kolmogorov complexity of the two branches can be bounded independently of the Born weights (for example, one branch is a short computational history and the other is an incompressible dump).

\paragraph{3. Holographic, not field-theoretic, vacuum energy.}
Zero-point sums are not performed.
The residual $\Lambda$ is a short program's output, and the UV cutoff of matter loops in a region is the area bound~\eqref{eq:holo}, not a cutoff in momenta of a continuum of modes.
This is already the standard holographic expectation~\cite{tHooft1993,bao2017}; here it is a theorem of the merge.

\paragraph{4. Moderate inflation, simple potentials.}
The measure $W$ rewards inflaton plugins that are short programs producing enough e-folds to solve the horizon problem, and penalizes both no inflation (long data on a large slice) and baroque landscapes named pointwise.
A single-field monomial or exponential of low degree is favoured over a random $10^{500}$-vacuum list unless the list itself has a short generator.
This is a selection principle on the landscape, not a unique potential.

\paragraph{5. Macroscopic solutions are Kolmogorov-elite.}
The observed catalogue---Kerr, FLRW, gravitational waves as linearized Einstein radiation, no naked singularities---is the catalogue of short GR-machine programs.
A stable, macroscopic, non-Einstein, non-diffeomorphic-to-Einstein geometry would falsify Einsteinian dominance (Proposition~\ref{prop:einstein-dom}) and so falsify the theory.

\paragraph{6. No physical intermediate cardinals, no non-separable kinematics.}
A compelling empirical role for a non-separable Hilbert space, or for a physical observable whose spectrum has cardinality strictly between $\alef$ and $\cont$, would falsify Postulates~\ref{post:qft}--\ref{post:sample}.
(Continuous spectra of ordinary operators are not a counterexample: they are $\cont$-sized, and they arise as inflations of countable spectral data, as in the spectral theorem for $L^{2}(\R)$.)

\paragraph{7. Black-hole entropy is program length.}
Any microscopic counting of black-hole states that does not, at leading order, reproduce~\eqref{eq:KBH} as a Kolmogorov bound is counting something other than the weight the merge uses.
Agreement of string, loop, and Euclidean calculations with $A/4\ell_{P}^{2}$ is, on this view, agreement that those calculations are enumerating shortest GR-machine descriptions.

\paragraph{8. No typical halt oracle.}
A $W$-typical, finite-area observer does not receive a clean halt bit from an infinite computation in finite proper time (Proposition~\ref{prop:no-halt}).
A laboratory or astrophysical Malament--Hogarth protocol that worked would falsify Einsteinian dominance or the holographic bound.

\section{Related work}

The theory sits at a crossing that several programs have approached from one side.

\paragraph{Algorithmic cosmologies.}
Solomonoff's universal prior~\cite{solomonoff1964}, Levin's semimeasure~\cite{levin1974}, and Schmidhuber's algorithmic theories of everything~\cite{schmidhuber2000} put $2^{-K}$ on universe-histories.
They do not put QFT and GR at two Cantorian infinities, they do not inflate samples with the Einstein Cauchy problem, and they do not merge.
M\"uller's algorithmic idealism~\cite{mueller2017,mueller2024} starts from observer states and Solomonoff prediction; the present theory starts from $\ket{\Omega}$ and a GR-machine, and treats observers as typical samples.

\paragraph{Complexity in quantum gravity.}
Dzhunushaliev's identification of gravitational probability with $e^{-K[g]}$, his rewriting of the path integral, and his Schwarzschild estimate~\cite{dzhunushaliev1998,dzhunushaliev2002} are the nearest ancestors of Postulate~\ref{post:weight}.
The present theory adds the cardinal split, the sampling of a separable $\ket{\Omega}$, the three merges, and the identification of holography with a Kraft bound on $\KG$.
Nielsen (circuit) complexity and Krylov complexity in cosmology~\cite{bhattacharyya2020,nielsen2006} are different quantities: they measure depth of a unitary, not program length of a universe.
They may bound $\KG$ in special states; they do not replace it.

\paragraph{Quantum cosmology.}
The Hartle--Hawking no-boundary state~\cite{hartle1983}, Vilenkin's tunneling state~\cite{vilenkin1983}, and decoherent histories~\cite{gellmann1993,hartle2008} are the coherent and decoherent merges in a different language, with $I_{E}$ in place of $\KG$.
The present theory claims that $I_{E}$ is the continuum shadow of $\KG$ on saddles, and that the sum is over programs, not over a continuum of metrics.

\paragraph{Relativistic hypercomputation.}
Malament--Hogarth spacetimes and Kerr-based relativistic computers~\cite{earman1993,hogarth2004,etesi2002,nemeti2006,welch2008} show that classical GR can decide HALT.
Section~\ref{sec:hyper} grants the classical construction and denies that $W$ samples it.

\paragraph{Holography and finite-dimensional gravity.}
't~Hooft, Susskind, Maldacena~\cite{tHooft1993,susskind1995,maldacena1998}; Bousso's bound~\cite{bousso1999}; Bao, Carroll, and Singh's locally finite-dimensional Hilbert space~\cite{bao2017}.
Equation~\eqref{eq:holo} is that circle of ideas stated as a bound on gravitational Kolmogorov complexity of a sample.

\paragraph{Discrete quantum gravity.}
Causal sets~\cite{bombelli1987}, CDT~\cite{ambjorn2005}, loop quantum gravity~\cite{rovelli2004,fairbairn2004}, and spin foams are candidate integer skeletons and candidate inflation maps.
The present theory is agnostic among them so long as $\Gamma$ is an Einstein completion and $\Hil$ is separable.

\paragraph{Cantorian spacetime.}
El Naschie's $E$-infinity Cantorian spacetime~\cite{elnaschie2004} uses transfinite Cantor sets as a model of spacetime and is a different project: it replaces the manifold, rather than treating the manifold as an inflation of an integer quantum skeleton.

\section{A discrete analog}

The $k$-cycle can be run in miniature, with no claim that the miniature is gravity.

Let $\Hil=\C^{2^{n}}$ with the computational basis $\{\ket{s}:s\in\{0,1\}^{n}\}$ as integer skeleton.
Let $\ket{\Omega}$ be a state of low circuit complexity (a short-program state).
Sampling is a computational-basis measurement, producing a bitstring $s$ of length $n$.
Inflation $\Gamma$ is, in the analog, the discrete Poisson problem---the nearest cousin of a constraint equation---or simply the identification of $s$ with a classical spin configuration on a cycle of length $n$.
Gravitational Kolmogorov complexity is approximated by
\begin{equation}
  \widehat{\KG}(s) \;=\; \ell_{\mathrm{gz}}(s) \;+\; \beta\, E_{\mathrm{smooth}}(s),
\end{equation}
where $\ell_{\mathrm{gz}}$ is the compressed length of $s$ (a computable upper bound on $K$) and $E_{\mathrm{smooth}}=\sum_{i}(s_{i+1}-s_{i})^{2}$ is a discrete curvature penalty, the analog of $I_{\mathrm{off}}$.
Weights $W(s)\propto\abs{\braket{s}{\Omega}}^{2}\,2^{-\widehat{\KG}(s)}$ are formed, and the merge is the weighted barycentre $\bar{s}=\sum_{s}W(s)\,s$ in $[0,1]^{n}$, or the typical sample.

Three things happen, and they are the three things the continuum theory claims.
Low-complexity, smooth strings dominate.
The merge is a finite sum.
The output is again a vector in a separable space (here finite-dimensional).
A companion program in this repository runs that analog.

The analog is also a warning.
Gzip is not $K$, a cycle graph is not Einstein, and $n<\infty$ erases the cardinal split.
What survives is the shape of the cycle: sample from a quantum state, score by compression plus a geometric penalty, merge.

\subsection{Minisuperspace analog}

A second companion program makes $\Gamma$ an Einstein solver.
The skeleton is finite-bit Cauchy data $(a,H,\phi,\dot\phi)$ for flat or closed FLRW with a homogeneous scalar (a cosmological constant, or $\tfrac12 m^{2}\phi^{2}$).
The on-shell slice solves the Friedmann constraint for $H$; those samples do not pay bits for Hubble, because Einstein supplies it.
Generic $4$-data specify $H$ independently and are charged
\begin{equation}
  I_{\mathrm{off}} \;=\; \bigl|3\bigl(H^{2}+k/a^{2}\bigr)-\rho\bigr|.
\end{equation}
Inflation is RK4 integration of the minisuperspace Einstein--Klein--Gordon system.
The weight is the same $k$-weight,
$\widehat{\KG}(s)=K(\mathrm{data}\mid\mathrm{Einstein})+\beta I_{\mathrm{off}}^{2}$.
On a $3$-bit grid the merge moves mass from the on-shell slice up relative to the Born measure, the typical universe is an expanding Einsteinian de~Sitter, and $\langle a(t)\rangle$ grows exponentially.
The algebraic core---Kraft's inequality for prefix-free programs, Friedmann constraint, de~Sitter as a fixed point of the minisuperspace vector field, Einstein as a bit-saving compressor, and $2^{\beta}$ dominance of on-shell vacuum over off-shell data---is a Lean 4 module in the repository.
Minisuperspace is not $3{+}1$ gravity, but it is the smallest place the compressor claim can fail.

The Kolmogorov vacuum is no longer a Gaussian. Born weights are the leading WKB minisuperspace amplitudes, in units $8\pi G=1$:
\begin{equation}
  |\Psi_{\mathrm{HH}}|^{2} \;\propto\; e^{+24\pi^{2} I_{0}}, \qquad
  |\Psi_{\mathrm{T}}|^{2} \;\propto\; e^{-24\pi^{2} I_{0}},
\end{equation}
with Halliwell's under-barrier integral $I_{0}$ equal to $1/V$ at the closed turning point $a^{2}=3/V$ (and equal to $1/V$ for $k=0$ nucleation).
On a quadratic inflaton the merge then does the textbook split: Hartle--Hawking concentrates on small $V$ and few e-folds; Vilenkin concentrates on large $V$ and more e-folds.
That is the no-boundary versus tunneling contest as a $\KG$ comparison on a finite skeleton. The remaining vacuum problem is to lift this from minisuperspace to $3{+}1$.

\subsection{Evaporation analog}

A third companion program makes $\Gamma$ a discrete Vaidya law in $1{+}1$.
The skeleton is an $N$-bit microstate $s$ of a hole of area $N$ (units $4\ell_{P}^{2}=1$, so area $=$ bits $=$ Bondi mass), a retarded time $t$, and a remaining mass $m$.
Each tick converts one Planck unit of mass into outgoing flux (constant $1{+}1$ luminosity).
On-shell, Einstein supplies the remaining mass: $m+t=N$, and the off-shell cost is $I_{E}=\abs{m+t-N}$.
Holography of the leftover hole is $I_{H}=\max(0,\text{bits in the hole}-m)$.
Three families of programs for the same $s$ are merged:
\begin{itemize}
  \item \emph{Page:} radiation is a reversible function of $s$; leftover bits equal $m$; $\widehat{\KG}$ stays $K(s)$ along the history~\cite{page1993}.
  \item \emph{Remnant:} all $N$ bits stay inside as $m$ drops; $I_{H}=t$; $\widehat{\KG}=K(s)+\beta t^{2}$.
  \item \emph{Scramble:} independent incompressible radiation of length $t$, interior still holding $s$; $\widehat{\KG}=K(s)+t+\beta t^{2}$.
\end{itemize}
The weight is the $k$-weight,
$\widehat{\KG}=K(\mathrm{data}\mid\mathrm{Einstein})+\beta(I_{E}^{2}+I_{H}^{2})$.
On an $N=6$ grid the merge is democratic at $t=0$ and Page-dominated by $t=N/2$; remnants and scrambles are exponentially downweighted; Schwarzschild ($s=0^{N}$) is the typical Page universe.
The algebraic core---Bondi constraint, leftover holography, constant Page $\widehat{\KG}$, and $2^{\beta t^{2}}$ dominance of Page over remnants---is a Lean 4 module in the repository.
Kraft's inequality is the static bound; this analog is that bound along a history.
$1{+}1$ Vaidya is not $3{+}1$ evaporation, but it is the smallest place the information-problem reading of~\eqref{eq:KBH} can fail.

\subsection{Laboratory Born-correction analog}

A fourth companion program runs~\eqref{eq:born-corr} on an $n$-qubit register, with no Einstein solver: $\Gamma$ is the identity, the sample is the laboratory record, and $\widehat{\KG}=K(\mathrm{record})$ in a tiny prefix language.
Two branches have independently bounded $K$: a short history $0^{n}$ ($K=2$) versus an incompressible dump ($K=n+1$).
Three scorings are compared.
Isolated, $\widehat{\KG}=K(\mathrm{register})$: equal Born on $n=6$ gives $\Delta K=5$ and merge mass $W_{\mathrm{short}}=0.97$ against Born $1/2$, the ratio $2^{\Delta K}$ of~\eqref{eq:born-corr}.
Two named clicks $0^{n}$ versus $1^{n}$ have $\Delta K=0$ even when isolated, so $W$ tracks Born---ordinary detector outcomes are both short programs.
A thermal bath of length $L$ is scored as $\widehat{\KG}=\max\bigl(K(\mathrm{env}),K(\mathrm{register})\bigr)$: once $L\ge K(\mathrm{dump})$ both records cost $L$ and $W$ tracks Born again.
Additive scoring $K(\mathrm{env})+K(\mathrm{register})$ is the control: $\Delta K$ is independent of $L$, so a large fridge does not by itself cancel an isolated-register dump.
The algebraic core---isolated $2^{c}$ bias, equal-$K$ clicks, additive non-swamp, thermal swamp at $K(\mathrm{env})\ge K(\mathrm{register})$---is a Lean 4 module in the repository.
A qubit register is not a universe, and the tiny language is not $K$.
What survives is the isolation theorem the paper warned about:~\eqref{eq:born-corr} is accessible only in a fully specified, isolated device whose total $\KG$ is dominated by the register, and it is not a claim that laboratory Born frequencies fail.

\subsection{Dimensionality analog}

A fifth companion program makes Dzhunushaliev's complexity-driven dimensional reduction~\cite{dzhunushaliev1998,dzhunushaliev2002} a merge on a finite catalog of gravity plugins $(D,L,\text{moduli})$.
Local graviton polarizations are $D(D-3)/2$.
The $p$-th Lovelock density~\cite{lovelock1971} is identically zero for $D<2p$, topological at $D=2p$, and dynamical for $D>2p$: Einstein--Hilbert is $p=1$; Gauss--Bonnet is topological in $4$d and dynamical for $D\ge 5$.
Plugins that do not compress Cauchy data---zero local gravitons, or no dynamical term---pay $I_{\mathrm{dead}}$.
$\widehat{\KG}=K(\mathrm{plugin})+\beta I_{\mathrm{dead}}$, with $4$d Einstein--Hilbert the $2$-bit vacuum; other $D$ pay $\abs{D-4}$, extra Lovelock terms and moduli their bits.
On the catalog $D=2{\ldots}6$ with optional Gauss--Bonnet and moduli in $\{0,1,2,4\}$, the typical plugin is $4$d Einstein--Hilbert, the $D=4$ family takes merge mass $0.55$ against $D=5$ at $0.30$ and $2{+}1$ at $0.001$, and compressors take $0.997$ against Born $0.40$.
Choquet--Bruhat well-posedness for $D\ge 3$ remains an input.
The algebraic core---$\mathrm{dof}(4)=2$, $\mathrm{dof}(3)=0$, Gauss--Bonnet topological in $4$d, and $2^{1+\beta}$ dominance of $4$d Einstein over dead $2{+}1$---is a Lean 4 module in the repository.
A catalog is not $3{+}1$ gravity.
What survives is the compressor count: $2{+}1$ is dead, extra dimensions and curvature junk are longer, $3{+}1$ Einstein is Kolmogorov-elite.

\subsection{Plugin language}

The dimensionality analog still listed gravity plugins by hand.
The Standard Model problem is the same encoding with matter.
A sixth companion program is that language, not a new $\Gamma$.
A group word is a 2-bit tagged prefix-free code
\begin{equation}
  G \;\coloneqq\; \mathrm{U}(1) \;\mid\; \mathrm{SU}(n) \;\mid\; \mathrm{SO}(n) \;\mid\; G\times H,
\end{equation}
with nats unary ($1^{n}0$).
A plugin is gravity (the dimensionality word) plus $G$, a fermion-rep count, $n_{\mathrm{gen}}$, and $\deg V$.
A raw landscape of $N$ vacua is $N$ concatenated copies, length $\Theta(N)$; a generator that prints $N$ copies costs $O(\log N)$.
Conditioned on a low-energy sample that requires $4$d Einstein and SM quantum numbers---an input, like Choquet--Bruhat---$SU(5)$ plus a 2-bit breaking program is length $24$ and takes merge mass $0.94$; the SM product is length $34$; a raw list of $16$ SM copies is length $308$.
$SU(5)$ is a \emph{shorter} word than $SU(3)\times SU(2)\times U(1)$.
The language does not derive the SM uniquely.
What it derives is that ``the SM is a plugin'' is a minimization over a fixed grammar, that GUTs win as short generators unless the breaking sector is charged in the word, and that a named vacuum list of size $N$ is a long program.
The algebraic core---$|\mathrm{SM}|=17$, $|\mathrm{SU}(5)|=8$, raw $N\times\mathrm{U}(1)$ longer than the SM word for $N\ge 5$, a $16$-copy generator shorter than the raw list, Kraft on the named encodings---is a Lean 4 module in the repository.

\section{Open problems}

\begin{enumerate}
  \item \textbf{The Einstein primitive.}
  Specify the GR-machine in enough detail that $L_{\mathrm{Einstein}}$ is a definite (if enormous) integer, and prove a precise version of~\eqref{eq:action-K} for Euclidean saddles.
  A minisuperspace analog (FLRW plus a scalar, Friedmann as compressor) is in the repository; the remaining task is a definite $L_{\mathrm{Einstein}}$ in $3{+}1$.

  \item \textbf{Uniqueness of $\Gamma$.}
  Different discretizations of Cauchy data give different maps $\Gamma$.
  Prove that their merges agree on all diffeomorphism-invariant coarse observables in the $\hbar\to 0$, $N\to\infty$ limit.

  \item \textbf{The Kolmogorov vacuum.}
  Construct, or at least uniquely characterize, $\ket{\Omega}$ as a shortest diffeomorphism-invariant state.
  The minisuperspace analog now runs Hartle--Hawking against tunneling as WKB amplitudes; the remaining task is the same contest in $3{+}1$.

  \item \textbf{Standard Model as plugin.}
  A finite plugin language is in the repository: prefix-free group words, $SU(5)$ shorter than the SM product, raw $N$-landscapes longer than either.
  The remaining task is to charge the GUT Higgs/breaking sector honestly enough that the minimizer on the grammar, given the low-energy sample, either \emph{is} the SM or fails in public.

  \item \textbf{Evaporation.}
  A $1{+}1$ Vaidya analog is in the repository: Page $\widehat{\KG}$ of hole plus radiation stays $K(s)$ along the history, while remnants and scrambles pay $\beta t^{2}$.
  The remaining task is the same conservation in $3{+}1$, including whether the matching to~\eqref{eq:KBH} is exact (no firewall, no remnant).

  \item \textbf{Complexity corrections.}
  A few-qubit interferometer analog is in the repository: isolated short-versus-dump disagrees with Born by $2^{\Delta K}$; thermal records with $L\ge K(\mathrm{dump})$, and two short clicks, track Born; additive $K(\mathrm{env})+K(\mathrm{reg})$ does not swamp.
  The remaining task is a bound in a concrete device, including whether $\Delta K$ can be made larger than laboratory systematics.

  \item \textbf{Dimensionality.}
  A catalog analog is in the repository, and gravity is now a nonterminal of the plugin language rather than a hand list.
  The remaining task is to prove that $3{+}1$ Einstein is the shortest well-posed compressor \emph{in that grammar}, not merely among $D=2{\ldots}6$ with a Gauss--Bonnet bit.
\end{enumerate}

\section{The theory in one page}

Cantor split physics into two infinities before physics knew it needed them.
The integers index quantum states.
The reals index spacetime.
The reals are the power set of the integers, so a universe is a subset of the quantum skeleton, completed by Einstein's equation.

Sample a finite string from $\ket{\Omega}$.
Inflate it with GR to a continuum geometry.
Weight it by $2^{-\KG}$ times its Born factor, where $\KG$ is Kolmogorov complexity on a machine that already knows Einstein.
Merge the weighted cloud.
The merge is a countable sum.
Classical spacetime is the typical term.
Quantum field theory is the integer residue.
Holography is Kraft's inequality on a region of finite area.
Black-hole entropy is program length.
The arrow of time is the unfolding of a short program.
The gravitational path integral was never an integral over the continuum.
It was always this sum.
Malament--Hogarth hypercomputation is a long program, not a typical merge.

That is a theory of everything in the only sense the phrase can bear.
It names the ontology, the measure, and the maps between the two infinities the twentieth century actually used.
It recovers GR as Einsteinian dominance of the merge, and QFT as residual integer dynamics.
It leaves the Standard Model as a shortest-plugin problem, $\KG$ as uncomputable, and $\ket{\Omega}$ as a shortest vector rather than a constructed one.
Those are debts.
They are debts of the right shape.

\begin{thebibliography}{99}

\bibitem{ambjorn2005}
J.~Ambj{\o}rn, J.~Jurkiewicz, and R.~Loll,
``Reconstructing the universe,''
Phys.\ Rev.\ D \textbf{72}, 064014 (2005).

\bibitem{bao2017}
N.~Bao, S.~M.~Carroll, and A.~Singh,
``The Hilbert space of quantum gravity is locally finite-dimensional,''
Int.\ J.\ Mod.\ Phys.\ D \textbf{26}, 1743013 (2017),
arXiv:1704.00066.

\bibitem{bekenstein1973}
J.~D.~Bekenstein,
``Black holes and entropy,''
Phys.\ Rev.\ D \textbf{7}, 2333 (1973).

\bibitem{bekenstein1981}
J.~D.~Bekenstein,
``Universal upper bound on the entropy-to-energy ratio for bounded systems,''
Phys.\ Rev.\ D \textbf{23}, 287 (1981).

\bibitem{bhattacharyya2020}
A.~Bhattacharyya, S.~Das, S.~Shajidul Haque, and B.~Underwood,
``Cosmological complexity,''
Phys.\ Rev.\ D \textbf{101}, 106020 (2020).

\bibitem{birrell1982}
N.~D.~Birrell and P.~C.~W.~Davies,
\textit{Quantum Fields in Curved Space}
(Cambridge University Press, Cambridge, 1982).

\bibitem{bombelli1987}
L.~Bombelli, J.~Lee, D.~Meyer, and R.~D.~Sorkin,
``Space-time as a causal set,''
Phys.\ Rev.\ Lett.\ \textbf{59}, 521 (1987).

\bibitem{bousso1999}
R.~Bousso,
``A covariant entropy conjecture,''
JHEP \textbf{07}, 004 (1999).

\bibitem{cantor1874}
G.~Cantor,
``\"Uber eine Eigenschaft des Inbegriffes aller reellen algebraischen Zahlen,''
J.\ Reine Angew.\ Math.\ \textbf{77}, 258 (1874).

\bibitem{cantor1878}
G.~Cantor,
``Ein Beitrag zur Mannigfaltigkeitslehre,''
J.\ Reine Angew.\ Math.\ \textbf{84}, 242 (1878).

\bibitem{cantor1891}
G.~Cantor,
``\"Uber eine elementare Frage der Mannigfaltigkeitslehre,''
Jahresber.\ Dtsch.\ Math.-Ver.\ \textbf{1}, 75 (1891).

\bibitem{carroll2017}
S.~M.~Carroll,
``Why Boltzmann brains are bad,''
arXiv:1702.00850.

\bibitem{chaitin1969}
G.~J.~Chaitin,
``On the length of programs for computing finite binary sequences: statistical considerations,''
J.\ ACM \textbf{16}, 145 (1969).

\bibitem{chaitin1974}
G.~J.~Chaitin,
``Information-theoretic limitations of formal systems,''
J.\ ACM \textbf{21}, 403 (1974).

\bibitem{choquetbruhat1952}
Y.~Foures-Bruhat,
``Th\'eor\`eme d'existence pour certains syst\`emes d'\'equations aux d\'eriv\'ees partielles non lin\'eaires,''
Acta Math.\ \textbf{88}, 141 (1952).

\bibitem{cohen1963}
P.~J.~Cohen,
``The independence of the continuum hypothesis,''
Proc.\ Natl.\ Acad.\ Sci.\ USA \textbf{50}, 1143 (1963).

\bibitem{dzhunushaliev1998}
V.~D.~Dzhunushaliev,
``Kolmogorov's algorithmic complexity and its probability interpretation in quantum gravity,''
Class.\ Quantum Grav.\ \textbf{15}, 603 (1998),
arXiv:gr-qc/9512014.

\bibitem{dzhunushaliev2002}
V.~Dzhunushaliev and D.~Singleton,
``Algorithmic complexity in cosmology and quantum gravity,''
Entropy \textbf{4}, 3 (2002),
arXiv:gr-qc/0108038.

\bibitem{earman1993}
J.~Earman and J.~D.~Norton,
``Forever is a day: supertasks in Pitowsky and Malament--Hogarth spacetimes,''
Philos.\ Sci.\ \textbf{60}, 22 (1993).

\bibitem{elnaschie2004}
M.~S.~El Naschie,
``A review of $E$-infinity theory and the mass spectrum of high energy particle physics,''
Chaos Solitons Fractals \textbf{19}, 209 (2004).

\bibitem{etesi2002}
G.~Etesi and I.~N\'emeti,
``Non-Turing computations via Malament--Hogarth space-times,''
Int.\ J.\ Theor.\ Phys.\ \textbf{41}, 341 (2002),
arXiv:gr-qc/0104023.

\bibitem{everett1957}
H.~Everett, III,
`` `Relative state' formulation of quantum mechanics,''
Rev.\ Mod.\ Phys.\ \textbf{29}, 454 (1957).

\bibitem{fairbairn2004}
W.~Fairbairn and C.~Rovelli,
``Separable Hilbert space in loop quantum gravity,''
J.\ Math.\ Phys.\ \textbf{45}, 2802 (2004),
arXiv:gr-qc/0403047.

\bibitem{faizal2025}
M.~Faizal, L.~Krauss, A.~Shabir, and F.~Marino,
``Consequences of undecidability in physics on the theory of everything,''
arXiv:2507.22950.

\bibitem{gellmann1993}
M.~Gell-Mann and J.~B.~Hartle,
``Classical equations for quantum systems,''
Phys.\ Rev.\ D \textbf{47}, 3345 (1993).

\bibitem{gibbons1977}
G.~W.~Gibbons and S.~W.~Hawking,
``Action integrals and partition functions in quantum gravity,''
Phys.\ Rev.\ D \textbf{15}, 2752 (1977).

\bibitem{gibbons1979}
G.~W.~Gibbons, S.~W.~Hawking, and M.~J.~Perry,
``Path integrals and the indefiniteness of the gravitational action,''
Nucl.\ Phys.\ B \textbf{138}, 141 (1979).

\bibitem{godel1940}
K.~G\"odel,
\textit{The Consistency of the Continuum Hypothesis}
(Princeton University Press, Princeton, 1940).

\bibitem{halmos1974}
P.~R.~Halmos,
\textit{Naive Set Theory}
(Springer, New York, 1974).

\bibitem{halliwell1991}
J.~J.~Halliwell and J.~B.~Hartle,
``Wave functions constructed from an invariant sum over histories satisfy constraints,''
Phys.\ Rev.\ D \textbf{43}, 1170 (1991).

\bibitem{hartle1983}
J.~B.~Hartle and S.~W.~Hawking,
``Wave function of the universe,''
Phys.\ Rev.\ D \textbf{28}, 2960 (1983).

\bibitem{hartle2008}
J.~B.~Hartle, S.~W.~Hawking, and T.~Hertog,
``The classical universes of the no-boundary quantum state,''
Phys.\ Rev.\ D \textbf{77}, 023502 (2008).

\bibitem{hawking1973}
S.~W.~Hawking and G.~F.~R.~Ellis,
\textit{The Large Scale Structure of Space-Time}
(Cambridge University Press, Cambridge, 1973).

\bibitem{hawking1975}
S.~W.~Hawking,
``Particle creation by black holes,''
Commun.\ Math.\ Phys.\ \textbf{43}, 199 (1975).

\bibitem{hogarth2004}
M.~Hogarth,
``Deciding arithmetic using SAD computers,''
Brit.\ J.\ Philos.\ Sci.\ \textbf{55}, 681 (2004).

\bibitem{kay1991}
B.~S.~Kay and R.~M.~Wald,
``Theorems on the uniqueness and thermal properties of stationary, nonsingular, quasifree states on spacetimes with a bifurcate Killing horizon,''
Phys.\ Rep.\ \textbf{207}, 49 (1991).

\bibitem{kolmogorov1965}
A.~N.~Kolmogorov,
``Three approaches to the quantitative definition of information,''
Probl.\ Inf.\ Transm.\ \textbf{1}, 1 (1965).

\bibitem{levin1974}
L.~A.~Levin,
``Laws of information conservation (nongrowth) and aspects of the foundation of probability theory,''
Probl.\ Inf.\ Transm.\ \textbf{10}, 206 (1974).

\bibitem{li2008}
M.~Li and P.~Vit\'anyi,
\textit{An Introduction to Kolmogorov Complexity and Its Applications}, 3rd~ed.
(Springer, New York, 2008).

\bibitem{lovelock1971}
D.~Lovelock,
``The Einstein tensor and its generalizations,''
J.\ Math.\ Phys.\ \textbf{12}, 498 (1971).

\bibitem{maldacena1998}
J.~Maldacena,
``The large $N$ limit of superconformal field theories and supergravity,''
Adv.\ Theor.\ Math.\ Phys.\ \textbf{2}, 231 (1998).

\bibitem{mueller2017}
M.~P.~M\"uller,
``Law without law: from observer states to physics via algorithmic information theory,''
arXiv:1712.01826.

\bibitem{mueller2024}
M.~P.~M\"uller,
``Algorithmic idealism: what should you believe to experience next?,''
arXiv:2412.02826.

\bibitem{mukhanov1992}
V.~F.~Mukhanov, H.~A.~Feldman, and R.~H.~Brandenberger,
``Theory of cosmological perturbations,''
Phys.\ Rep.\ \textbf{215}, 203 (1992).

\bibitem{nemeti2006}
I.~N\'emeti and G.~D\'avid,
``Relativistic computers and the Turing barrier,''
Appl.\ Math.\ Comput.\ \textbf{178}, 118 (2006).

\bibitem{nielsen2006}
M.~A.~Nielsen,
``A geometric approach to quantum circuit complexity,''
Science \textbf{311}, 1133 (2006); quant-ph/0502070.

\bibitem{page1993}
D.~N.~Page,
``Information in black hole radiation,''
Phys.\ Rev.\ Lett.\ \textbf{71}, 3743 (1993).

\bibitem{rovelli2004}
C.~Rovelli,
\textit{Quantum Gravity}
(Cambridge University Press, Cambridge, 2004).

\bibitem{schmidhuber2000}
J.~Schmidhuber,
``Algorithmic theories of everything,''
arXiv:quant-ph/0011122.

\bibitem{solomonoff1964}
R.~J.~Solomonoff,
``A formal theory of inductive inference, I and II,''
Inf.\ Control \textbf{7}, 1 (1964); \textbf{7}, 224 (1964).

\bibitem{susskind1995}
L.~Susskind,
``The world as a hologram,''
J.\ Math.\ Phys.\ \textbf{36}, 6377 (1995).

\bibitem{tHooft1993}
G.~'t~Hooft,
``Dimensional reduction in quantum gravity,''
in \textit{Salamfestschrift},
ed.\ A.~Ali, J.~Ellis, and S.~Randjbar-Daemi
(World Scientific, Singapore, 1993),
arXiv:gr-qc/9310026.

\bibitem{vilenkin1983}
A.~Vilenkin,
``Birth of inflationary universes,''
Phys.\ Rev.\ D \textbf{27}, 2848 (1983).

\bibitem{welch2008}
P.~D.~Welch,
``The extent of computation in Malament--Hogarth spacetimes,''
Brit.\ J.\ Philos.\ Sci.\ \textbf{59}, 659 (2008),
arXiv:gr-qc/0609035.

\bibitem{wald1994}
R.~M.~Wald,
\textit{Quantum Field Theory in Curved Spacetime and Black Hole Thermodynamics}
(University of Chicago Press, Chicago, 1994).

\bibitem{wolfram2020}
S.~Wolfram,
\textit{A Project to Find the Fundamental Theory of Physics}
(Wolfram Media, Champaign, 2020).

\bibitem{zurek1989}
W.~H.~Zurek,
``Algorithmic randomness and physical entropy,''
Phys.\ Rev.\ A \textbf{40}, 4731 (1989).

\end{thebibliography}

\end{document}
